%% kapa: exam=1 round=33 year=2022 period=1 form=A subject=civil_law subject_ko=민법 seq=39 number=39
%% source: Q-Net 감정평가사 기출문제 (cst00309) artlSeq=5209997 file=raw/1008/5209997/2022년도 제33회 감정평가사 1차시험 1교시(민법,경제학,부동산학).pdf
%% conversion: kordoc 4.24.51 -> tools/gen1.py (auto); review=proofread-sonnet; flags=none
\begin{question}{39}
근저당권에 관한 설명으로 옳지 않은 것은? (다툼이 있으면 판례에 따름)
\begin{choices}
  \item 근저당권의 존속기간이나 결산기를 정한 경우, 원칙적으로 결산기가 도래하거나 존속기간이 만료한 때에 그 피담보채무가 확정된다.
  \item 근저당권의 존속기간이나 결산기를 정하지 않고 피담보채권의 확정방법에 관한 다른 약정이 없는 경우, 근저당권설정자는 근저당권자를 상대로 언제든지 계약 해지의 의사표시를 하여 피담보채무를 확정시킬 수 있다.
  \item 근저당권자가 피담보채무의 불이행을 이유로 경매신청을 한 경우, 경매신청시에 근저당권의 피담보채권액이 확정된다.
  \item 후순위 근저당권자가 경매를 신청한 경우, 선순위 근저당권의 피담보채권은 매수인이 매각대금을 완납한 때에 확정된다.
  \item 공동근저당권자가 저당목적 부동산 중 일부 부동산에 대하여 제3자가 신청한 경매절차에 소극적으로 참가하여 우선배당을 받은 경우, 특별한 사정이 없는 한 나머지 저당목적 부동산에 관한 근저당권의 피담보채권도 확정된다.
\end{choices}
\end{question}
\answer{5}
